\chapter{Assumptions and Promises}
\label{ch:assumptions}
\label{ch:generated-examples}

\index{assumption}
\index{promise}
A protocol says what \emph{can} happen. Two more things decide whether a system
is good enough: how often each outcome happens, and how long and how much it
costs. LeanFM keeps them apart. A probability on a choice is an
\emph{assumption}: what a role is expected to decide where nobody else can see.
A \texttt{require} clause is a \emph{promise}: a minimum or maximum the system
keeps, judged over an observation period. LeanFM checks the promises against the
assumptions before any code exists, and against the event stream once it runs.

\section{Assumptions on choices}

\index{choice!probability}
A choice whose decider's reasons are hidden carries the probabilities it is
assumed to have, and a name:

\begin{lstlisting}[language={},basicstyle=\ttfamily\footnotesize,breaklines=true]
import LeanFM.Session.Syntax

/-- A server that is sometimes busy, with one hard deadline and one promised fraction. -/
protocol PingServer (Client, Server) {
  q: Ping from Client to Server;
  choice at Server using declared("IsBusy") {
    probability 9/10:
      Pong from Server to Client cause q after 50 ms;
  } or {
    probability 1/10:
      SlowPong from Server to Client cause q after 800 ms;
  }
}

require per q {
  every q reaches Pong, SlowPong within 1 second;
  at least 95/100 of q reach Pong, SlowPong within 100 ms over 1000 entries;
}
\end{lstlisting}


\texttt{choice at Server using declared("IsBusy")} says the server decides, that
the probabilities are declared (stated by the author, not measured), and that the
decision is called \texttt{IsBusy}. The name belongs to the whole choice, like a
question; the branches are its answers. Reports name the answer by the first
message of the branch taken, so the same choice reads
\texttt{[IsBusy → Pong: 9/10]} on one branch and
\texttt{[IsBusy → SlowPong: 1/10]} on the other. Probabilities are exact
fractions and must sum to one. An overlay can replace them by name without
editing the protocol, for example with rates measured from a log
(Chapter~\ref{ch:bakery}).

\section{Time, amounts and rates}

\index{delay}
\index{amount}
\index{rate}
Messages can take time and move value. None of these annotations changes the
protocol's structure: projection, conformance and temporal questions are the
same with or without them.

\begin{center}
\small
\begin{tabular}{@{}p{.47\textwidth}p{.46\textwidth}@{}}
\toprule
Annotation & Meaning \\
\midrule
\texttt{after 2 days}, \texttt{after sample \{ 1 days : 1/2; 3 days : 1/2; \}} & a delay after the previous step, fixed or drawn from a distribution \\
\texttt{amount 1000 USD}, \texttt{amount sample \{ \ldots \}} & value moved from sender to receiver \\
\texttt{amount \{ 5 apple; -3 USD; \}} & a trade vector: positive entries move forward, negative back \\
\texttt{accrue 10 USD per week from A to B;} & a rate in force until the next \texttt{accrue} for the same roles and unit \\
\texttt{at next 1 week} & a periodic gate: wait for the next whole week since the session began \\
\texttt{on breach record Late pays 2000 USD to T;} & in an obligation: the penalty when its deadline passes \\
\bottomrule
\end{tabular}
\end{center}

Units are names, so \texttt{roomweek} (one room for one week) is as good a unit as
\texttt{USD}. Results are kept per role and per unit.

\index{decision tree}
\index{Markov chain}
With its assumptions, a protocol without loops is a decision tree: every path to
an end has an exact probability, a time, and each role's balance. LeanFM draws
it:

\begin{figure}[htbp]
\centering
\AccessibleFigure{Decision tree of Ping. Start. Ping from Client to Server at 0 ms. Server chooses IsBusy. With policy 9/10 Pong from Server to Client at 50 ms ends with terminal probability 9/10. With policy 1/10 SlowPong at 800 ms ends with terminal probability 1/10.}{%
\ThemedDiagram{height=.72\textheight}{lang-ping-tree}}
\caption{\texttt{Ping} under its assumption, as LeanFM draws it: the server's
choice with its probabilities, each answer's time, and each ending's exact
probability.}
\label{fig:ping-tree}
\end{figure}

A protocol with loops (\texttt{rec}) is solved instead as an absorbing Markov
chain. Expected visits $v$ solve $v = s + vQ$ over the transient states, exactly,
in rational arithmetic, so a retry loop that runs a geometric number of times
has an exact expected cost rather than a truncated one. A session that some state
can never leave is reported as one that can loop forever.

\section{Promises}

\index{require}
\index{service level}
Promises are written after the protocols, one block per entry message. Every
occurrence of the entry is one trial. \texttt{Ping}'s block says every ping is
answered within a second, and at least 95 in 100 within 100 ms, judged over 1000
pings.

\begin{center}
\small
\begin{tabular}{@{}p{.42\textwidth}p{.51\textwidth}@{}}
\toprule
Clause & Meaning \\
\midrule
\texttt{every q reaches A, B within d;} & a hard deadline: each \texttt{q} reaches \texttt{A} or \texttt{B} within \texttt{d}, and one miss is a counterexample \\
\texttt{at least p of q reach A within d over W;} & a promised minimum fraction \\
\texttt{at most p of q reach A over W;} & a promised maximum fraction \\
\texttt{over 1 hour}, \texttt{over 500 entries} & the observation period, required for every fraction \\
\texttt{confidence 99/100} & how sure a verdict must be (default 95/100) \\
\texttt{expecting 500 per hour} & the expected traffic, so a time window's size can be checked \\
\bottomrule
\end{tabular}
\end{center}

\index{observation period}
A fraction needs an observation period. Without one, the first slow answer after
deployment would break a promise about ninety-five answers in a hundred before
there was anything to judge it by. A window is judged only once it has closed.

LeanFM checks each promise against the assumptions on the chain, following every
path from each entry with exact probabilities until a target is reached in time
(a hit) or the deadline, the next entry or the end comes first (a miss):

\begin{lstlisting}[language={},basicstyle=\ttfamily\footnotesize,breaklines=true]
Requirements against the stated assumptions (exact, per entry):
- every q reaches Pong, SlowPong within 1 second;
  MET: per entry the probability of reaching Pong, SlowPong within 1 second is 1.
- at least 95/100 of q reach Pong, SlowPong within 100 ms over 1000 entries;
  VIOLATED: per entry the probability of reaching Pong, SlowPong within 100 ms is 9/10. The most probable path that breaks it: Ping -> [IsBusy -> SlowPong: 1/10] -> SlowPong -> (at 800 ms, past the 100 ms deadline), through the assumption [IsBusy -> SlowPong: 1/10]. The window holds the 60 entries it needs to decide at 19/20 confidence. The assumptions cannot meet this bound: change the protocol (a retry, a fallback) or the stated probabilities on the path shown.
\end{lstlisting}


The deadline holds: even the slow answer comes at 800 ms. The fraction fails,
because only 9 in 10 answers come within 100 ms, and the report names the
assumption responsible. The fix is a faster busy path or a weaker promise. A
window needs $\lceil \ln(1/(1-c))/m \rceil$ entries to tell the promise apart
from a failure rate of 0 or 1, where $m$ is the margin to the bound; LeanFM
reports when a window cannot hold that many.

A larger example, a user lookup through a proxy that can lose a forward and a
database that can fail:

\lstinputlisting[language={},basicstyle=\ttfamily\scriptsize,breaklines=true]{../examples/sessions/UserLookup.leanfm}

\lstinputlisting[language={},basicstyle=\ttfamily\scriptsize,breaklines=true]{generated/sla-userlookup.txt}

The error budget cannot be met: $\frac{999}{1000}\cdot\frac{1}{100} +
\frac{1}{1000} = \frac{1099}{100000}$ is more than $\frac{1}{200}$, and the report
names both assumptions on the path. Meeting it is a design change, such as a
retry after \texttt{Failure}. The window needs 600 entries and an hour at 500 an
hour holds only 500, so the promise needs a longer window or more traffic before
it can be judged at all.
